\section{Applicability beyond Singapore}
\label{sec:applicability}

The encoding presented in \S\ref{sec:evaluation} is Singapore-only,
and the paper is careful in \S\ref{sec:limitations} not to overstate
portability without empirical evidence. This section makes the
\emph{positive} case: where the grammar plausibly transfers, what
adaptation each candidate jurisdiction would require, and what
empirical work would falsify the claim. We separate this from the
limitations discussion deliberately. Disclaiming a thing and arguing
for its plausibility are different rhetorical moves, and
proof-of-concept readers reasonably want both.

\subsection{The Anglo-Indian penal code family}
\label{subsec:anglo_indian_family}

The Singapore Penal Code 1871 is not a Singapore invention. It is the
Straits Settlements adoption of the 1860 Indian Penal Code drafted by
Macaulay's First Law Commission, lightly localised and subsequently
amended. The same parent text underpins the criminal codes of India
(IPC~1860, since 2023 reframed as the Bharatiya Nyaya Sanhita),
Pakistan (PPC~1860), Bangladesh (PC~1860), Brunei (PC, Cap.~22),
Malaysia (PC, Act~574), Sri Lanka (PC No.~2 of 1883), Myanmar
(PC~1861), and historically several East African codes. The shared
ancestry is not merely thematic. Section numbering, marginal-note
conventions, illustration formatting, exception clauses, and the
characteristic actus-reus / mens-rea / circumstance partitioning are
recognisably the same drafting style across the family. A reader who
has read Singapore's s415 will find the corresponding Indian and
Malaysian s415 cosmetically familiar, often differing only in
penalty quanta and amendment lineage.

This is what makes the family the natural test bed for grammar
portability. Yuho's primitives were chosen to fit the drafting shape
of penal statute --- not the cultural specifics of one jurisdiction.
The element partitioning is doctrinal, not local; the exception block
is a Catala-style prioritised override pattern that maps directly
onto the codified ``provided that'' clauses common to every member of
the family; the penalty combinators capture the cumulative-or-
alternative-or-conditional structure that the entire family inherits
from Macaulay's draft. To the extent the grammar fits Singapore
because the grammar fits Macaulay, the grammar should fit the others
for the same reason.

\subsection{What transfers, what adapts}
\label{subsec:portability_mechanism}

We separate the grammar substrate (which we expect to transfer
unchanged) from the per-jurisdiction configuration (which would need
swapping).

\paragraph{Transfers as-is.} The tree-sitter grammar, the AST, the
five-block statute structure (\texttt{elements}, \texttt{penalty},
\texttt{exceptions}, \texttt{illustrations}, \texttt{caselaw}), the
G10 reference resolver, the Z3 / Alloy verification scaffolding, the
seven transpilers (including the OASIS Akoma Ntoso target), and the entire editor surface (LSP, MCP, browser
extension, explorer site, Word add-in) are jurisdiction-neutral. They
operate on the AST, not on Singapore-specific tokens. Adopting a new
jurisdiction would not require parser changes for these subsystems.

\paragraph{Adapts via configuration.} A small number of surface
features are configurable per jurisdiction:

\begin{itemize}
\item \emph{Penalty vocabulary.} The set of admissible punishments
varies: caning is in the Singapore, Malaysian, and Bruneian codes
but not the Indian one; Pakistan retains \emph{tazir} categories
absent from Singapore; the Bangladeshi 2020 amendments enumerate
specific aggravators absent elsewhere. Yuho's punishment kinds are
a closed enum in the grammar; the per-jurisdiction adoption needs
the enum extended (or restricted) and the controlled-English
transpiler taught the new vocabulary.
\item \emph{Currency tokens.} Fines are encoded as a numeric quantum
plus a currency tag (\texttt{SGD} for Singapore). Indian, Pakistani,
and Bruneian codes use different currencies; the type system already
admits the parameter, but the L2 fidelity linter expects the
jurisdiction's reporting currency on initial scrape.
\item \emph{Statutory amendment lineage.} The \texttt{amends} clause
points at an act number (e.g.\ ``Act 32 of 2019''); each
jurisdiction has its own amendment-act numbering scheme. The grammar
accepts arbitrary strings here; only the L3 fidelity check that
reconciles \texttt{amends} against the canonical scrape needs a
jurisdiction-specific source.
\item \emph{Marginal-note style.} Older codes (1860--1900) place the
marginal note above the section; newer ones (post-1990) inline it.
The Phase~A scraper's heuristic for extracting marginal notes is
jurisdiction-specific; the AST and downstream tooling are not.
\end{itemize}

\paragraph{Requires substantive work.} Two surfaces are not yet
configuration-clean:

\begin{itemize}
\item \emph{Mens-rea taxonomy.} Singapore distinguishes intention,
knowledge, rashness, and negligence as distinct mental states; the
1860 IPC is identical, but the BNS~2023 (the recent Indian
recodification) introduces collapsed categories that do not map
one-to-one. A jurisdiction that uses a different taxonomy needs
either a grammar extension or a vocabulary remapping pass.
\item \emph{The L3 reviewer's checklist.} The 11-point fidelity
checklist (\S\ref{subsec:coverage}) bakes in conventions specific to
the Singapore SSO format (anchor IDs, marginal-note syntax,
amendment notation). A jurisdiction with a different official
publication format would need its own checklist; the
parameterisation surface exists but has not been exercised.
\end{itemize}

\subsection{Beyond the Penal Code family}
\label{subsec:non_penal}

The grammar is statute-shaped rather than penal-statute-shaped at the
top level: a module contains statutes, a statute contains elements
and penalties and exceptions, and the AST does not require the
penalties block to be present. The question of whether Yuho fits
non-penal statutes (tax, employment, intellectual property,
competition law) reduces to whether those statutes share the
\emph{element + exception + illustration} shape. The Singapore
Income Tax Act and the Employment Act both exhibit this shape on
inspection of a sample of sections; the Companies Act does too,
though with a heavier reliance on definitions blocks that Yuho
already supports. A more rigorous claim awaits empirical encoding.

The shift to civil-law jurisdictions is more speculative. The French
\emph{Code p\'enal} and the German \emph{Strafgesetzbuch} are
similarly element-structured, but their definitional density and
their reliance on cross-references to the civil-law general parts
mean an encoding would lean heavily on Yuho's reference graph and
less on its element partitioning. We do not claim the grammar fits
civil-law codes; we claim the AST is the right level at which to
test the question. Catala~\cite{catala2021}, focused on French tax
law, is the closer analogue here, and a Yuho-Catala interop layer
(scope composition over Yuho statute blocks) is the most natural
extension if the civil-law direction is pursued.

\subsection{An empirical path to validation}
\label{subsec:portability_path}

We commit to a concrete and falsifiable test of the portability
claim. Encoding a 50-section sample of the Indian Penal Code is the
minimum viable exercise: 50 sections is enough to surface grammar
gaps that the Singapore exercise did not, and small enough to
complete in a single encoder-quarter. The sampling strategy mirrors
the Singapore encoding order: start with structurally simple
property offences (theft, cheating, criminal misappropriation),
move through bodily-harm hierarchies (hurt, grievous hurt, culpable
homicide), and end with the procedural and abetment apparatus.
Concretely, the test passes if (a)~the grammar absorbs the 50-section
sample at L1 with zero new gaps catalogued, (b)~the L2 fidelity
linters produce diagnostic densities within an order of magnitude of
the Singapore baseline, and (c)~the existing G1--G14 gap catalogue
covers any drafting patterns the IPC sample exhibits but Singapore's
PC does not. Failure on (a) implies a fifteenth gap and a grammar
extension; failure on (b) implies the linters are over-fitted to
Singapore conventions; failure on (c) implies the gap catalogue is
not a stable artefact and would need a replication study before any
portability claim is publishable.

The corpus directory layout is already jurisdiction-aware: the slim
corpus, the canonical per-section JSON, and the reference graph are
keyed on a \texttt{jurisdiction} field that defaults to Singapore but
admits any string. Switching corpus base from Singapore to India
requires no code change to the editor surfaces, only a parallel
\texttt{library/penal\_code\_in/} directory and a per-jurisdiction
\texttt{config.toml} for the per-jurisdiction parameters above. The
build pipeline already supports running the corpus build over an
arbitrary library root, so a multi-jurisdiction artefact is an
engineering exercise on top of the empirical encoding work.

We frame this as future work in \S\ref{sec:conclusion} rather than as
a present-tense claim. What this section adds beyond
\S\ref{sec:limitations} is a positive structural argument: the
grammar substrate is jurisdiction-neutral, the configuration surface
is small, and the empirical test is bounded. Portability has not been
demonstrated, but the cost of demonstrating it is well-understood and
the path to falsification is clear.
